% CECI EST UN FICHIER DE MODÈLE LATEX POUR LES ARTICLES INCLUS DANS
% *Anthology of Computers and the Humanities*.
% 'final' LORS DE LA CRÉATION DE LA VERSION FINALE DE L'ARTICLE.
% NE PAS modifier la classe du document
%\documentclass[final,fra]{anthology-ch} % pour la version finale
\documentclass[final,fra]{anthology-ch}         % pour la soumission

% CHARGER LES PACKAGES LaTeX
\usepackage{booktabs}
\usepackage{graphicx}

% AJOUTEZ vos propres packages avec \usepackage{}

% TITRE DE LA SOUMISSION
% Remplacez ceci par le titre de votre soumission
\title{Modéliser une tradition de fables ésopiques : architecture TEI, taxonomies et structuration des connaissances dans l'édition numérique de l'Isopet 1-Avionnet}

% INFORMATIONS SUR LES AUTEURS ET LES AFFILIATIONS
% Pour chaque auteur, utilisez une nouvelle commande \author, avec
% les numéros entre crochets indiquant les affiliations correspondantes
% (section suivante) et l'ORCID-ID pour chaque auteur.
\author[1]{Joana Casenave}[
  orcid=0000-0002-1609-1512,
  email=joana.casenave@univ-lille.fr
]

% Il doit y avoir un appel à \affiliation pour chaque affiliation des auteurs.
% Plusieurs affiliations peuvent être attribuées à un auteur
% et une affiliation peut être attribuée à plusieurs auteurs.
\affiliation{1}{Université de Lille, France}

% MOTS-CLÉS
% Fournissez un ou plusieurs mots-clés séparés par des virgules
% en utilisant les commandes suivantes
\keywords[fra]{édition critique numérique, XML/TEI, modélisation des données, taxonomies et thésaurus, annotation, structuration des connaissances, fables médiévales}

\keywords[eng]{critical digital edition, XML/TEI, data modeling, taxonomies and thesaurus, annotation, knowledge structure, medieval fables}
%\keywords[eng]{informatique, sciences humaines}

% MÉTADONNÉES DE LA PUBLICATION
% Ces champs seront remplis lors de la publication ;
% ils peuvent être laissés par défaut lors de la soumission
\pubyear{2026}
\pubvolume{4}
\pagestart{211}
\pageend{219}
\paperorder{26}
\conferencename{Actes de la Conférence Humanistica}
\conferenceeditors{Serena Crespi, Simon Gabay, Martin Grandjean, Ariane Pinche, Marie Puren et Léa Saint-Raymond}
\doi{10.63744/oFjQfi7IwgQB}
\customhead{}

\addbibresource{bibliography.bib}

%%%%%%%%%%%%%%%%%%%%%%%%%%%%%%%%%%%%%%%%%%%%%%%%%%%%%%%%%%%%%%%%%%%%%%%%%%%
% DÉBUT DU TEXTE
\begin{document}

\maketitle

\begin{abstract}
This article examines knowledge modeling within the framework of digital scholarly editing, drawing on the corpus of medieval fables from the Isopet 1-Avionnet as part of the FabuLab project. It proposes a stratified TEI architecture that distinguishes between a conceptual level - comprising a master TEI file that consolidates taxonomies and reference frameworks - and a documentary level corresponding to the encoding of manuscript witnesses. This structure makes it possible to separate the definition of categories from their use in texts, while ensuring their articulation through mechanisms of reference and annotation. The study shows that this organization promotes consistency in annotations, reduces redundancy, and enables corpus exploration based on analytical categories. The article thus advocates a conception of digital scholarly editing as a space for the formalization of knowledge, in which analytical categories are integrated into the data model.
\end{abstract}


\section{Introduction}\label{introduction-1}

La conception d'éditions critiques numériques conduit à repenser les
modalités de modélisation et de représentation des connaissances
produites par le travail philologique.
Dans le cas des fables ésopiques
médiévales, la diversité des traditions textuelles, la richesse des
motifs narratifs et la multiplicité des catégories d'analyse
mobilisables font apparaître des enjeux spécifiques liés à la
construction de modèles de connaissances.
Dans cette perspective,
l'édition critique numérique est envisagée avant tout
comme un dispositif de structuration, de formalisation et d'exploitation
des connaissances à l'aide de la TEI.
Cette réflexion s'inscrit dans le
cadre du projet FabuLab, dont l'un des objectifs est de
concevoir une architecture d'encodage articulant édition des textes,
indexation savante et élaboration de systèmes de catégories et de
référentiels.

Le corpus étudié, l'\emph{Isopet 1-Avionnet}, présente une configuration
particulièrement propice à une réflexion sur ces questions : coexistence
de deux rédactions successives, variations structurelles importantes,
forte densité d'intertextualité et caractéristiques formelles et
génériques propres à la fable \cite{boivin2006naissance, casenave2023fables}.
Se pose dès
lors une double question: comment traduire ces caractéristiques dans un
modèle d'encodage TEI qui ne se limite pas à représenter les textes,
mais qui explicite les structures, les catégories et les relations qui
les organisent?
Et comment concevoir, à partir de ce cadre, des
instruments d'exploration du corpus fondés sur des taxonomies, des index
structurés et des thésaurus? L'enjeu est de développer un modèle capable
de rendre compte à la fois de la structure des textes, des dispositifs
d'annotation savante et des systèmes classificatoires mobilisés pour
l'analyse du corpus.
Le présent article propose d'examiner les principes
retenus dans ce cadre et d'analyser la manière dont l'encodage TEI peut
être mobilisé pour articuler édition critique, représentation des
connaissances et dispositifs d'indexation.
\section{État de l'art}\label{uxe9tat-de-lart}

La réflexion s'inscrit dans le champ désormais bien établi de l'édition
critique numérique, qui a montré que l'encodage ne saurait être réduit à
une opération de balisage technique, mais constitue une véritable
activité de modélisation des objets textuels et des hypothèses
interprétatives \cite{apollon2014digital, pierazzo2015digital, driscoll2016digital, sahle2016what, bleier2018digital,  casenave2023edition}.
Dans
le prolongement des travaux de Jerome McGann sur l'édition conçue comme
\emph{knowledge site}, de nombreux auteurs ont souligné que les modèles
de données jouent un rôle structurant dans l'organisation et
l'opérationnalisation des savoirs \cite{pierazzo2019what}.
L'édition numérique
savante peut ainsi être comprise comme un dispositif intellectuel dans
lequel le savoir critique n'est pas seulement formulé dans le discours,
mais est aussi inscrit dans l'architecture même des données, des catégories et des
relations qui les organisent.

Dans cette perspective, la modélisation apparaît comme une pratique
centrale des humanités numériques : non plus simple opération de
représentation, mais véritable processus de production de connaissances
fondé sur la construction et la mobilisation de structures formelles.
Arianna Ciula \cite{ciula2017digital, ciula2023modelling} propose de l'envisager
en tant que stratégie de recherche permettant de produire et
d'interpréter des représentations externes afin de
donner sens à des objets complexes.
Cette conception conduit à
considérer les dispositifs d'encodage comme des espaces
de médiation entre données textuelles et structures conceptuelles, au
sein desquels se construisent et se stabilisent les catégories de
l'analyse.
Elle invite également à penser la textualité
non plus en objet stable, mais en objet dynamique, défini par des
relations de production, de transmission et de réécriture.

Parallèlement, on observe ces dernières années une évolution marquée
vers des modèles d'édition intégrant de plus en plus explicitement des
dispositifs d'alignement sémantique et d'interopérabilité.
Les
taxonomies locales en TEI tendent ainsi à être alignées sur des
vocabulaires contrôlés et des modèles de type SKOS \cite{baker2013key},
tandis que l'alignement des données devient un enjeu central, inscrivant
progressivement l'édition savante dans l'écosystème des graphes de
connaissances \cite{ciotti2014tei, eide2014ontologies}.
Plusieurs travaux ont
montré que la TEI peut être comprise comme une forme de modélisation
conceptuelle propre aux humanités numériques, distincte mais
complémentaire des ontologies formelles utilisées dans les systèmes
d'information patrimoniaux \cite{ciula2014reflections}.
Des recherches ont
ainsi exploré le recours à des ontologies RDF/OWL comme extension des
encodages TEI, afin de représenter non seulement des entités mais aussi
les relations complexes qui les relient et d'expliciter des associations
implicites dans les sources \cite{ciula2008expressing}.
Dans le domaine des
éditions savantes numériques, ces réflexions ont conduit à explorer
l'intégration des données textuelles encodées avec des modèles de
graphes et des technologies du web sémantique, notamment afin de
représenter et d'exploiter les relations complexes entre textes,
témoins, acteurs et concepts \cite{spadini2021graph}.

Toutefois, malgré ces avancées, la formalisation explicite des systèmes
conceptuels internes aux éditions numériques demeure relativement peu
interrogée.
Si les éditions mobilisent fréquemment des catégories
critiques, génériques ou intertextuelles pour structurer les corpus, ces
catégories restent souvent implicites dans les modèles de données ou
confinées à des taxonomies locales peu formalisées.
La question de leur
modélisation - c'est-à-dire de la manière dont ces catégories sont
construites, structurées et articulées au sein des données éditoriales -
reste encore largement sous-théorisée.
Or, dans le cas de corpus
fortement structurés par des classifications génériques ou des relations
intertextuelles, la formalisation de ces systèmes catégoriels devient un
enjeu central, tant pour l'intelligibilité scientifique des éditions que
pour leur interopérabilité et leur réutilisation dans des environnements
de recherche élargis.

Dans cette perspective, le présent article propose d'examiner la manière
dont un système de catégories critiques peut être explicitement modélisé
au sein d'une édition numérique, en considérant le document TEI non
seulement comme un format d'encodage, mais aussi comme une structure de
modélisation capable d'organiser et de formaliser un référentiel
conceptuel interne au corpus.
L'analyse prend pour cas d'étude le corpus
de l'\emph{Isopet 1-Avionnet} et s'appuie sur une architecture TEI
permettant de rendre explicite l'articulation entre données textuelles
et structures conceptuelles (voir Figure 1).

\section{Le corpus : l'\emph{Isopet 1-Avionnet}}

L'\emph{Isopet 1-Avionnet} occupe une position singulière dans la
tradition médiévale des fables ésopiques.
Issu de la convergence des
deux principales branches occidentales de cette tradition, le
\emph{Romulus} et l'\emph{Avianus}, il se présente comme une compilation
fondée sur plusieurs recueils sources, l'\emph{Anonyme de Nevelet} pour
l'\emph{Isopet 1} et l'\emph{Avianus} pour l'\emph{Avionnet}.
Le texte a
par ailleurs connu deux états rédactionnels successifs : une première
version, puis une réécriture opérée quelques décennies plus tard par un
compilateur, qui introduit des ajouts, des réorganisations et des
transformations substantielles \cite{boivin2006naissance}.

La tradition manuscrite de l'\emph{Isopet 1-Avionnet} se compose de six
témoins, désignés par les sigles établis par K. McKenzie et W. Oldfather \cite{McKenzie_ysopet_1919}, et répartis de manière équilibrée entre les deux états du texte.
Les manuscrits \textit{B}\footnote{Manuscrit \textit{B} : Bruxelles, KBR, ms. 11193, 134~f.}, \textit{P}\footnote{Manuscrit \textit{P} : Paris, BnF, ms. fr.~1594, f. 1r-113v.}
et \textit{L}\footnote{Manuscrit \textit{L} : Londres, BL, Add. 33781, 139 f.}, datés du \textsc{xiv}\textsuperscript{e} siècle, transmettent la seconde rédaction, tandis
que le groupe
\textit{a}\footnote{Manuscrit \textit{a} : Paris, BnF, ms.fr. 1595, f. 1r-36r.}
\textit{b}\footnote{Manuscrit b : Paris, BnF, ms. fr. 19123, f.~110r-132v.}, 
et \textit{c}\footnote{Manuscrit \textit{c} : Paris, BnF, ms. fr.~24310, f.~2r-57v.}, 
daté du \textsc{xv}\textsuperscript{e} siècle, conserve la première version, antérieure.

Les témoins de la seconde rédaction présentent en outre la particularité
d'intégrer directement les sources latines, placées en amont de chaque
fable en moyen français.
À ces sources explicites s'ajoute un réseau
dense de références intertextuelles, qui déborde le cadre strict des
recueils de fables pour inclure d'autres œuvres médiévales ainsi que des
éléments du répertoire parémiologique.

Cette stratification des sources, des états du texte et des réseaux de
références fait de l'\emph{Isopet 1-Avionnet} un terrain
particulièrement favorable à une approche fondée sur la modélisation des
connaissances.
Elle soulève notamment la question de la manière dont ces
différentes dimensions peuvent être représentées et articulées dans un
dispositif éditorial cohérent.
Le projet FabuLab se présente comme la
préparation d'une édition critique numérique exhaustive de ce corpus, en
mobilisant les potentialités des technologies d'encodage pour structurer
les différentes couches du texte et proposer des parcours de lecture
multiples fondés sur l'organisation même des données.

\section{Méthodologie}\label{muxe9thodologie}

Pour répondre à cet enjeu, le projet repose sur une architecture TEI
stratifiée visant à articuler deux niveaux complémentaires au sein d'une
même infrastructure (Figure 1) :
\begin{enumerate}
    \item un niveau conceptuel, correspondant aux systèmes de catégories et aux référentiels mobilisés pour l'analyse du corpus ;
    \item un niveau documentaire, correspondant aux textes des témoins manuscrits et à leur encodage.
\end{enumerate}

Cette architecture s'organise autour d'un fichier TEI maître, qui
constitue le socle conceptuel commun, et d'un ensemble de fichiers TEI
correspondant chacun à un témoin manuscrit.
Ce découplage répond à une
double exigence : éviter la redondance des annotations analytiques dans
les différents témoins, tout en garantissant la cohérence et la
stabilité des catégories utilisées.

\begin{figure}[htp]
\centering
\includegraphics[width=1.0\textwidth]{./media/image2.png}
\caption{Architecture à deux niveaux du modèle TEI pour la modélisation des textes et des connaissances (projet FabuLab).}
\end{figure}

Le fichier maître regroupe l'ensemble des référentiels
partagés du corpus et les systèmes de catégorisation : taxonomies
analytiques (types de fables, structures discursives, formes de
moralités, actes de discours), ainsi que des index décrivant les entités
mentionnées dans les textes (personnages des fables, lieux concrets /
abstraits, objets, référents zoologiques, notions abstraites, entités
religieuses et mythologiques, types humains et rôles sociaux,
événements, personnages historiques, œuvres référencées dans le texte).
Ces ressources sont structurées à l'aide des éléments
standards de la TEI (notamment \texttt{<classDecl>} pour les taxonomies et
\texttt{<listPerson>}, \texttt{<listPlace>}, \texttt{<listObject>}, \texttt{<listNym>}, \texttt{<listEvent>},\texttt{<listBibl>} pour les
index).
Elles sont toutes associées à des identifiants pérennes qui
permettent leur mobilisation dans les différents fichiers textuels.
Ce
choix permet de considérer ces catégories non comme de simples aides à
l'annotation, mais comme des objets de données autonomes, réutilisables
et interopérables.
Le fichier maître accueille également la
documentation scientifique globale du corpus et constitue l'espace
commun des métadonnées éditoriales.

Les fichiers TEI correspondant aux témoins manuscrits sont encodés selon un schéma homogène et reliés au fichier maître par des mécanismes de référence. Les entités mentionnées dans les textes sont annotées à l'aide de l'élément \texttt{<rs>}, dont l'attribut \texttt{@ref} pointe vers les identifiants du fichier maître, tandis que les catégories analytiques des taxonomies sont mobilisées via des attributs de typage (notamment \texttt{@ana}). Ce dispositif permet de dissocier la définition des catégories de leur mobilisation dans les textes : il maintient une séparation claire entre modélisation et encodage, tout en garantissant leur articulation.

Afin d'illustrer concrètement la mobilisation, dans l'encodage des
textes, des catégories définies dans le TEI maître, l'extrait suivant
(Figure 2) présente le début de la première fable du recueil, tel qu'il
est encodé dans le fichier du manuscrit \emph{c}.
Les catégories
déclarées dans les taxonomies analytiques et lexicales, ainsi que les
référentiels d'entités (index), sont mobilisées au moyen des attributs
analytiques et des mécanismes de référence pour structurer le texte et
en permettre l'interrogation.

\begin{figure}[htp]
    \centering
    \includegraphics[width=0.9\textwidth]{./media/image3.png}
    \caption{Extrait d'encodage TEI d'une fable : annotation des entités et des catégories analytiques.}
\end{figure}

Par ailleurs, l'annotation analytique n'est pas uniformément distribuée
sur l'ensemble des témoins.
Nous faisons l'hypothèse qu'une annotation
centralisée sur un nombre restreint de témoins de référence permet
d'éviter les effets de redondance et de divergence.
En conséquence, les
dispositifs d'annotation analytique et d'indexation sont appliqués
uniquement aux témoins éditoriaux de référence retenus pour chaque
rédaction -- le manuscrit \textit{c} pour la première rédaction et le manuscrit \textit{B}
pour la seconde.
Les autres témoins sont encodés comme témoins alignés,
principalement destinés à la comparaison textuelle et à l'étude des
variantes.
Dans une deuxième phase à venir du projet FabuLab, une
édition critique sera proposée pour chacune des deux rédactions de
l'Isopet 1-Avionnet et les témoins éditoriaux de
référence deviendront alors les manuscrits de base de cette édition
critique numérique.

L'articulation entre les témoins repose sur un système d'identifiants attribués aux unités textuelles (fables et vers). Ce dispositif constitue un élément clé de l'architecture : il permet d'aligner les différentes versions sans dupliquer les annotations, et de maintenir un lien entre variation textuelle et structures analytiques.

Ainsi conçue, l'architecture TEI ne se limite pas à un dispositif d'encodage, mais met en œuvre une véritable modélisation des connaissances, dans laquelle les catégories analytiques sont explicitement définies, stabilisées et mobilisables à l'échelle de l'ensemble du corpus.

\section{Résultats et discussion}\label{ruxe9sultats-et-discussion}

La mise en œuvre de cette architecture a conduit à la constitution
effective d'un environnement éditorial structuré autour d'un fichier TEI
maître et de six fichiers TEI correspondant aux témoins manuscrits du
corpus.
Cette organisation permet d'articuler, dans la pratique de
l'encodage et de l'exploitation des données, les différents niveaux de
description définis dans la modélisation initiale.
Elle permet également
de valider plusieurs hypothèses méthodologiques relatives à
l'articulation entre encodage, annotation et structuration des
connaissances.
La Figure 3 met en évidence la circulation et
l'articulation des données, depuis leur modélisation
dans le fichier TEI maître jusqu'à leur mobilisation dans l'encodage et
leur exploitation dans l'analyse du corpus.

\begin{figure}[htp]
\centering
\includegraphics[width=1.0\textwidth]{./media/image4.png}
\caption{Du TEI maître à l'analyse : circulation des
données entre modélisation, encodage et exploitation du corpus (projet
FabuLab).}
\end{figure}

L'un des premiers résultats concerne la centralisation effective des référentiels et des systèmes de catégorisation dans le fichier TEI maître. Celui-ci fonctionne comme un espace de stabilisation conceptuelle, dans lequel sont définies les catégories analytiques et les entités mobilisées dans l'annotation. Cette centralisation assure la cohérence des annotations et facilite leur réutilisation dans l'ensemble des fichiers textuels. Elle rend également possible l'alignement de ces référentiels avec les thésaurus et bases de connaissances externes
Biblissima\footnote{Le thésaurus iconographique de Biblissima est un  vocabulaire contrôlé développé dans le cadre de l'infrastructure de  recherche Biblissima, dédié à la description et à l'indexation des  images et motifs dans les manuscrits médiévaux.
\url{https://data.biblissima.fr}.},
Thema\footnote{ThEMA (\textit{Thesaurus Exemplorum Medii Aevi}) est une base de données multimédia et multilingue consacrée aux récits exemplaires médiévaux, dotée d'un thésaurus pour leur indexation thématique. \url{https://thema.huma-num.fr}.} 
et Wikidata\footnote{\url{https://www.wikidata.org}.}, inscrivant ainsi les données du corpus dans un réseau plus large de ressources sémantiques et de données interopérables.

Du point de vue éditorial, l'architecture retenue permet d'articuler transcription exhaustive des témoins et annotation analytique contrôlée. Les manuscrits \emph{c} et \emph{B}, retenus comme témoins de référence pour chacune des deux rédactions, concentrent les dispositifs d'annotation savante, tandis que les autres manuscrits servent principalement de points de comparaison textuelle. Cette organisation montre qu'il est possible de dissocier annotation savante et couverture documentaire, sans perte d'information pour l'analyse.

Enfin, la formalisation des catégories analytiques comme objets de
données constitue un résultat structurant. Les catégories rhétoriques et lexicales ne relèvent plus seulement du commentaire critique, mais sont intégrées au modèle sous forme de structures explicitement définies et interrogeables. Cette évolution ouvre la possibilité d'explorer le corpus selon des entrées conceptuelles (types de morale, structures discursives, catégories d'entités) et de croiser ces dimensions avec les variations textuelles.

Ainsi, l'architecture mise en place permet de penser conjointement la
variabilité des traditions manuscrites et la stabilité des cadres
d'analyse.
Elle fournit un environnement éditorial dans
lequel la production de l'édition critique et la structuration des
connaissances sont étroitement articulées, ouvrant la voie à des formes
d'exploration et d'analyse du corpus qui dépassent les modèles
éditoriaux traditionnels.

\section{Perspectives}\label{perspectives}

Dans une perspective d'ouverture vers le web sémantique, les systèmes
taxonomiques élaborés dans ce cadre ont vocation à être externalisés et
reconfigurés sous la forme d'une ontologie légère, notamment en SKOS \cite{baker2013key}.
Cette évolution vise également à étendre la modélisation à la structure
divergente du corpus et aux différentes formes d'intertextualité qu'il
met en jeu, afin de les intégrer dans un ensemble cohérent de relations
formalisées.

Ce travail s'inscrit dans une conception de l'édition numérique comme
instrument de connaissance, dans laquelle l'encodage ne constitue pas
seulement un moyen de représentation, mais une opération de modélisation
permettant d'expliciter et de formaliser les catégories de l'analyse.
L'architecture stratifiée mise en place dans ce projet, distinguant un
niveau conceptuel (fichier TEI maître, taxonomies, référentiels) et un
niveau documentaire (encodage des témoins manuscrits), participe
pleinement de cette approche.
Elle vise à faire du corpus non seulement
un ensemble de textes édités, mais aussi un espace structuré de
formalisation et d'exploration des savoirs relatifs aux fables
médiévales.

\section*{Financement}

Le projet FabuLab (dir. Joana Casenave) est réalisé grâce au soutien de l'ÉquipEx Biblissima+ (financement 2024-2026), qui bénéficie d'une aide de l’État gérée par l'ANR au titre du Programme d'investissements d'avenir intégré à France 2030, portant la référence ANR-21-ESRE-0005.


\nocite{*}
\printbibliography

\end{document}